\begin{textsample}{NEG Dim 2 – 2017\_c5 – Score -3.00 – 2017\_c5/enade\_2017\_pedag\_30ae56952d.txt}  \label{ex:f2_neg_044}
CAPÍTULO 5 DISTRIBUIÇÃO DOS CONCEITOS 5.1 PANORAMA NACIONAL DA DISTRIBUIÇÃO DOS CONCEITOS O Conceito Enade é um indicador de qualidade que avalia os cursos por intermédio dos desempenhos dos estudantes a partir dos resultados do Enade.
Os valores possíveis vão de 1 ( pior situação ) a 5 ( melhor situação ).
Os cursos com apenas um ou sem nenhum concluinte participante não obtêm o Conceito Enade, ficando Sem Conceito ( SC ).
A Tabela 5.1 apresenta a quantidade e distribuição de cursos de Pedagogia participantes do Enade/2017, por faixa de conceito e Grande Região. \textbf{Enfatiza} -se, mais uma vez, que a diferença entre os cursos tabulados neste capítulo e no capítulo 2 corresponde aos cursos Sem Conceito ( SC ), ou seja àqueles, em princípio, sem alunos concluintes que houvessem participado da prova ou àqueles que tivessem tão somente um aluno concluinte18.
Observando -se os dados da Tabela 5.1, nota -se que, dos 1.212 cursos participantes, 516 ( 42,6 \% ) classificaram -se com conceito 3, o valor modal.
Esse foi também o conceito modal em quase todas as regiões : Nordeste ( 40,8 \% ), Sudeste ( 43,9 \% ), Sul ( 47,9 \% ) e Centro-Oeste ( 38,8 \% ).
Na região Norte, a moda foi o conceito 2, atribuído para 46,6 \% dos cursos da região.
O conceito 2 foi o segundo mais frequente em nível nacional ( 25,8 \%, correspondendo a 313 cursos ) e o conceito 4, o terceiro ( 22,4 \%, correspondendo a 271 cursos ).
Houve, ainda, 58 cursos ( 4,8 \% ) que receberam conceito 5 e 43 cursos ( 3,5 \% ) que receberam conceito 1.
Dos 1.212 cursos de Pedagogia, 11 ( 0,9 \% ) ficaram Sem Conceito ( SC ).
A região Norte participou com 88 cursos ou 7,3 \% do total nacional.
Desses, 41 cursos receberam o conceito 2, o que equivale a 46,6 \% do total regional, além de ser o conceito modal na região, como já comentado. 18 Estes cursos com somente um concluinte não foram considerados no capítulo 2. modal na região, como já comentado.
A 11 cursos ( 12,5 \% ) atribuiu -se conceito 1, a 28 cursos ( 31,8 \% ), conceito 3 e a sete cursos ( 8,0 \% ), conceito 4.
Um curso ( 1,1 \% ) recebeu conceito 5 e nenhum curso ficou Sem Conceito ( SC ).
A região Nordeste participou com 228 cursos ou 18,8 \% do total nacional.
Nessa região, 93 cursos ( 40,8 \% em termos regionais ) obtiveram conceito 3, o conceito modal para a região, como já comentado.
Os conceitos 1 e 2 foram atribuídos, respectivamente, a 17 e a 77 cursos ( 7,5 \% e 33,8 \%, respectivamente ).
Ainda nessa região, 32 cursos ( 14,0 \% ) foram avaliados com o conceito 4, e sete cursos ( 3,1 \% ), com conceito 5.
Dois cursos ( 0,9 \% ) ficaram Sem Conceito ( SC ).
Dos 556 cursos participantes da região Sudeste ( 45,9 \% do total nacional ), o conceito 3 foi o conceito modal, como já assinalado, tendo sido obtido por 244 cursos ( 43,9 \% ).
O conceito 1 foi atribuído a cinco cursos ( 0,9 \% ) e o conceito 2, a 129 cursos ( 23,2 \% ).
Receberam os conceitos 4 e 5, respectivamente, 135 cursos ( 24,3 \% ) e 38 cursos ( 6,8 \% ).
Nessa região, cinco cursos ( 0,9 \% ) ficaram Sem Conceito ( SC ).
Os 211 cursos da região Sul corresponderam a 17,4 \% do total nacional.
Como já indicado, houve predominância do conceito 3, atribuído a 47,9 \% dos cursos da região ( conceito modal ), ou seja, atribuído a 101 dos 211 cursos participantes na região Sul.
O conceito 1 foi atribuído a um curso ( 0,5 \% ), o conceito 2, a 23 cursos ( 10,9 \% ), o conceito 4, a 72 cursos ( 34,1 \% ) e o conceito 5, a 11 cursos ( 5,2 \% ).
Três cursos ( 1,4 \% ) ficaram Sem Conceito ( SC ) na região Sul.
Já dos 129 cursos participantes na região Centro-Oeste ( 10,6 \% do total nacional ), 50 receberam conceito 3, o conceito modal ( 38,8 \% ), como já fora destacado.
A nove cursos foi atribuído conceito 1 ( 7,0 \% ), e a 43 cursos ( 33,3 \% ), o conceito 2.
Os conceitos restantes, 4 e 5, foram atribuídos, respectivamente, a 25 cursos ( 19,4 \% ) e a um curso ( 0,8 \% ).
Nessa região, um curso ( 0,8 \% ) ficou Sem Conceito ( SC ).
O Gráfico 5.1 apresenta a distribuição cumulativa do Conceito Enade segundo Grande Região.
Num gráfico de distribuição cumulativa, em cada valor de abscissa ( eixo x ), o valor apresentado no eixo das ordenadas ( eixo y ) é o valor acumulado da variável.
Por exemplo, a região Norte ( linha azul ) apresenta em torno de 90 \% dos cursos com conceito 3 ou abaixo deste valor.
A região Nordeste, por outro lado, apresenta cerca de 80 \% dos cursos na mesma situação : com conceito 3 ou abaixo.
Nas regiões Norte ( linha azul ) e Centro-Oeste ( linha preta ), todos os cursos têm conceitos 4 ou abaixo disso, já que a poligonal cumulativa alcança 100 \% para o conceito 4.
Em linhas gerais, num gráfico de distribuição cumulativa, poligonais mais à esquerda representam Grandes Regiões com pior distribuição, e poligonais mais à direita, Grandes Regiões com melhores conceitos.
A região Norte ( linha azul ) apresenta os cursos com conceitos mais baixos, ao passo que as regiões Sudeste ( linha vermelha ) e Sul ( linha roxa ) apresentam as distribuições com valores maiores.
Cursos em IES situadas nas regiões Nordeste ( linha verde ) e Centro-Oeste ( linha preta ) apresentam uma situação intermediária. 5.2 CONCEITOS POR CATEGORIA ADMINISTRATIVA, POR MODALIDADE DE ENSINO E GRANDE REGIÃO A Tabela 5.2 apresenta a distribuição dos cursos de Pedagogia participantes do Enade/2017, por Categoria Administrativa e Modalidade de ensino, de acordo com os conceitos por eles alcançados, segundo as Grandes Regiões brasileiras.
Dos 1.212 cursos participantes, 345 ( 28,5 \% ) eram ministrados em Instituições Públicas e 867 ( 71,5 \% ), em Instituições Privadas.
Quanto à Modalidade de ensino, 96 cursos eram a Distância e 1.116, Presenciais.
De acordo com as informações da Tabela 5.2, em termos nacionais, entre Instituições Públicas e Privadas, observa -se que, dos 58 cursos avaliados com conceito 5, 21 eram oferecidos em IES Públicas e 37, em IES Privadas.
Dos 345 cursos participantes de IES Públicas, o conceito 3 foi o valor modal, atribuído a 148 cursos ( 42,9 \% ).
Entre os demais cursos de Instituições Públicas participantes, 18 ( 5,2 \% ) outros 67 cursos ( 19,4 \% ) obtiveram conceito 2, 87 cursos ( 25,2 \% ) receberam conceito 4 e, como já comentado, 21 cursos ( 6,1 \% ) receberam conceito 5.
Nesta categoria, quatro cursos ( 1,2 \% ) ficaram Sem Conceito ( SC ).
Na rede privada, o conceito modal também foi 3, atribuído a 368 cursos, o correspondente a 42,4 \% dos 867 cursos da categoria.
Entre os demais cursos participantes, 25 ( 2,9 \% ) receberam conceito 1, e 246 ( 28,4 \% ), conceito 2.
O conceito 4 foi atribuído a 184 cursos ( 21,2 \% ) e o conceito 5, a 37 cursos ( 4,3 \% ).
Nesta Categoria Administrativa, sete cursos ( 0,8 \% ) ficaram Sem Conceito ( SC ).
Dos 1.116 cursos participantes da modalidade presencial, o conceito 3 também foi o valor modal, atribuído a 482 cursos ( 43,2 \% ).
Entre os demais cursos nesta modalidade de ensino, 250 cursos ( 22,4 \% ) foram avaliados com conceito 4, 278 cursos ( 24,9 \% ) obtiveram conceito 2, 55 cursos ( 4,9 \% ) receberam conceito 5 e 40 cursos ( 3,6 \% ) receberam conceito 1.
Nesta modalidade 11 cursos ( 1,0 \% ) ficaram Sem Conceito ( SC ).
Entre os cursos a Distância, o conceito modal foi 2, atribuído a 35 cursos, o correspondente a 36,5 \% dos 96 cursos da modalidade a Distância.
Entre os demais cursos participantes, 34 ( 35,4 \% ) receberam conceito 3, 21 ( 21,9 \% ), conceito 4 e três ( 3,1 \% ), conceito 5.
Nessa Modalidade nenhum curso ficou Sem Conceito ( SC ) e três cursos ( 3,1 \% ) receberam conceito 1.
Na análise por região, observa -se que, na região Norte, as Instituições Privadas participaram com 43 cursos ( 48,9 \% do total regional ), dos quais nenhum curso ficou Sem Conceito ( SC ).
O conceito modal para as Instituições Privadas na região Norte foi 2, atribuído a 24 cursos, correspondendo a 55,8 \% dessa categoria na região.
Três cursos ( 7,0 \% ) foram avaliados com conceito 1, 15 cursos ( 34,9 \% ) receberam conceito 3, nenhum curso recebeu conceito 4 e um curso ( 2,3 \% ) obteve o conceito 5.
As Instituições Públicas participaram com 45 cursos na região Norte ( 51,1 \% do total regional ), dos quais 17 obtiveram conceito 2, o valor modal ( 37,8 \% ).
O conceito 3 foi atribuído a 13 cursos ( 28,9 \% ), e o conceito 1, a oito cursos ( 17,8 \% ).
O conceito 4 foi atribuído a sete cursos ( 15,6 \% ).
Nenhum curso ficou Sem Conceito ( SC ) ou com conceito 5.
Na região Norte 96,6 \% dos cursos oferecidos eram presenciais e o conceito modal foi 2, obtido por 38 dos 85 cursos presenciais ( 44,7 \% ).
Os demais cursos presenciais foram alocados nos conceitos 1 ( 11 cursos, 12,9 \% ), 3 ( 28 cursos, 32,9 \% ), 4 ( sete cursos, 8,2 \% ) e 5 ( um curso, 1,2 \% ).
Os três cursos à Distância da região Norte foram alocados no conceito 2.
Na região Nordeste, a rede Privada concentrou 115 dos 228 cursos participantes, o equivalente a 50,4 \% do total da região.
Dentre os cursos oferecidos por IES Privadas no Nordeste, 55 foram avaliados com conceito 2 ( 47,8 \% ), o conceito modal.
Dois cursos obtiveram conceito 5 ( 1,7 \% ), dez cursos receberam conceito 1 ( 8,7 \% ), 12 cursos ficaram com conceito 4 ( 10,4 \% ) e 36 cursos alcançaram o conceito 3 ( 31,3 \% ).
Nessa combinação de categoria e região, nenhum curso ficou Sem Conceito ( SC ).
As Instituições Públicas da região Nordeste participaram com 113 cursos ( 49,6 \% do total da região ).
Desses, 57 obtiveram conceito 3 ( 50,4 \% ), o valor modal, seguidos por outros 22 cursos, que receberam conceito 2 ( 19,5 \% ).
O conceito 4 foi atribuído a 20 cursos ( 17,7 \% ).
Dois cursos ( 1,8 \% ) ficaram Sem Conceito ( SC ).
Outros cinco cursos receberam conceito 5 ( 4,4 \% ) e sete ( 6,2 \% ) ficaram com conceito 1 nessa região.
A grande maioria dos cursos na região Nordeste foram presencias : 213 ( 93,4 \% ).
Para estes, a moda foi o conceito 3 com 88 cursos ( 41,3 \% ).
Os demais cursos presenciais da região foram alocados aos conceitos 1 ( 16 cursos, 7,5 \% ), 2 ( 70 cursos, 32,9 \% ), 4 ( 30 cursos, 14,1 \% ) e 5 ( sete cursos, 3,3 \% ).
Dois cursos presenciais ( 0,9 \% ) ficaram Sem Conceito ( SC ).
Os conceitos dos 15 cursos a Distância ( 6,6 \% ) ficaram assim \textbf{distribuídos} : um curso com conceito 1 ( 6,7 \% ), sete cursos com o conceito 2 ( 46,7 \% ), cinco cursos com o conceito 3 ( 33,3 \% ) e dois cursos com o conceito 4 ( 13,3 \% ).
Nenhum curso a Distância da região foi alocado no conceito 5 ou ficou Sem Conceito ( SC ).
Na região Sudeste, a proporção de cursos da rede privada, 85,6 \%, foi mais \textbf{elevada} do que a registrada nas demais regiões brasileiras, correspondendo a 476 dos 556 cursos participantes.
Nessa categoria e região, o conceito modal foi 3, atribuído a 217 cursos ( 45,6 \% ).
Quatro cursos ( 0,8 \% ) ficaram Sem Conceito ( SC ).
Os demais foram avaliados com conceito 1 ( quatro cursos, 0,8 \% ), conceito 2 ( 117 cursos, 24,6 \% ), conceito 4 ( 108 cursos, 22,7 \% ) e conceito 5 ( 26 cursos, 5,5 \% ).
Entre os 80 cursos oferecidos por Instituições Públicas na região Sudeste ( 14,4 \% ), o conceito modal ocorreu nos conceitos 3 e 4, atribuídos a 27 cursos, 33,8 \% cada.
Os demais cursos receberam os conceitos 2 e 5 ( ambos com 12 cursos, 15,0 \% ) e conceito 1 ( um curso, 1,3 \% ).
Na região Sudeste, nesta categoria, um curso ( 1,3 \% ) ficou Sem Conceito ( SC ).
Também na região Sudeste, a maioria dos cursos foram presencias : 509 ( 91,5 \% ).
Para estes, a moda foi o conceito 3 com 228 cursos ( 44,8 \% ).
Os demais cursos presenciais da região foram alocados aos conceitos : 1 ( cinco cursos, 1,0 \% ), 2 ( 115 cursos, 22,6 \% ), 4 ( 121 cursos, 23,8 \% ), e 5 ( 35 cursos, 6,9 \% ).
Cinco cursos ( 1,0 \% ) ficaram Sem Conceito ( SC ).
Para os 47 cursos a Distância ( 8,5 \% ), a moda também foi o conceito 3 ( 34,0 \% ).
Os demais cursos a Distância da região foram alocados aos conceitos 2 e 4 ( 14 cursos, 29,8 \% em cada ) e 5 ( três cursos, 6,4 \% ).
Nenhum curso a Distância da região foi alocado ao conceito 1 e nenhum ficou Sem Conceito ( SC ).
As Instituições Privadas concentraram 153 dos 211 cursos participantes da região Sul, 72,5 \% do total regional.
Desses, 73 ( 47,7 \% ) obtiveram conceito 3, o conceito modal.
Nesta combinação de Categoria Administrativa e Grande Região, os conceitos 2 e 4 foram atribuídos a 21 ( 13,7 \% ) e a 49 ( 32,0 \% ) cursos, respectivamente.
Oito cursos ( 5,2 \% ) receberam conceito 5, nenhum recebeu conceito 1 e dois ( 1,3 \% ) ficaram Sem Conceito ( SC ).
As Instituições Públicas na região Sul participaram com 58 cursos, 27,5 \% dos cursos da região.
O conceito modal foi 3, atribuído a 28 cursos ( 48,3 \% ).
Aos demais cursos foram atribuídos os conceitos 1 ( um curso, 1,7 \% ), 2 ( dois cursos, 3,4 \% ), 4 ( 23 cursos, 39,7 \% ) e 5 ( três cursos, 5,2 \% ).
Um curso ( 1,7 \% ) da região Sul oferecido por IES Pública ficou Sem Conceito ( SC ).
Também na região Sul, a maioria dos cursos foram Presenciais, correspondendo a 189 ( 89,6 \% ), por oposição aos 22 cursos a Distância ( 10,4 \% ).
Para os cursos presenciais, a moda foi o conceito 3 com 90 cursos ( 47,6 \% ).
Os demais cursos presenciais da região foram alocados aos conceitos 2 ( 16 cursos, 8,5 \% ), 4 ( 69 cursos, 36,5 \% ) e 5 ( 11 cursos, 5,8 \% ).
Nenhum curso presencial foi alocado ao conceito 1 e três ( 1,6 \% ) ficaram Sem Conceito ( SC ).
Para os 22 cursos a Distância, a moda foi o conceito 3 ( 50,0 \% ).
Outros sete cursos a Distância da região ( 31,8 \% ) foram alocados ao conceito 2, três receberam conceito 4 ( 13,6 \% ) e um foi alocado ao conceito 1 ( 4,5 \% ).
Nenhum curso ficou Sem Conceito ( SC ) ou recebeu conceito 5.
Na região Centro-Oeste, 80 dos 129 cursos participantes eram de Instituições Privadas ( 62,0 \% em termos regionais ).
Desses, 29 receberam conceito 2 ( 36,3 \% ), o conceito modal.
Entre os demais, oito receberam conceito 1 ( 10,0 \% ), 27 receberam conceito 3 ( 33,8 \% ), 15 receberam conceito 4 ( 18,8 \% ).
Um curso ( 1,3 \% ) ficou Sem Conceito ( SC ) e nenhum curso recebeu conceito 5.
Quanto aos 49 cursos oferecidos por Instituições Públicas na região Centro-Oeste ( 38,0 \% do total regional ), 23 cursos apresentaram conceito 3 ( 46,9 \% ), conceito modal.
Um ficou com conceito 1 ( 2,0 \% ) e outro com o conceito 5 ( 2,0 \% ).
Nenhum curso ficou Sem Conceito ( SC ).
Já o conceito 4 foi atribuído a dez cursos ( 20,4 \% ) e o conceito 2 a 14 cursos ( 28,6 \% ).
Nesta região, 120 cursos foram Presenciais ( 93,0 \% ) e nove, a Distância ( 7,0 \% ).
Para os cursos presenciais, a moda foi o conceito 3 com 48 cursos ( 40,0 \% ).
Os demais foram alocados aos conceitos 1 ( oito cursos, 6,7 \% ), 2 ( 39 cursos, 32,5 \% ), 4 ( 23 cursos, 19,2 \% ) e 5 ( um curso, 0,8 \% ).
Um curso presencial ( 0,8 \% ) ficou Sem Conceito ( SC ).
Dos nove cursos a Distância quatro foram alocados no conceito 2 ( 44,4 \% ) e os demais ficaram \textbf{distribuídos} nos conceitos 1 ( um curso, 11,1 \% ), 3 ( dois cursos, 22,2 \% ) e 4 ( dois cursos, 22,2 \% ).
Nenhum curso ficou Sem Conceito ( SC ) ou recebeu conceito 5.
O Gráfico 5.2 apresenta a distribuição do Conceito Enade segundo a Categoria Administrativa da IES.
Os cursos em IES Públicas ( linha azul ) apresentam uma poligonal mais à direita do que os em IES Privadas ( linha verde ) e, consequentemente, uma distribuição de Conceitos Enade com valores maiores. 5.3 CONCEITOS POR ORGANIZAÇÃO ACADÊMICA E GRANDE REGIÃO Na Tabela 5.3, encontra -se a distribuição dos conceitos atribuídos aos cursos participantes do Enade/2017 na Área de Pedagogia, por Organização Acadêmica, segundo as Grandes Regiões brasileiras.
Dos 1.212 cursos de Pedagogia participantes, 494 eram oferecidos em Universidades, 141 em Centros Universitários, 570 em Faculdades e sete em CEFET/IFET.
Esta distribuição corresponde a, respectivamente, 40,8 \%, 11,6 \%, 47,0 \% e 0,6 \% dos cursos.
De acordo com os dados apresentados, de todos os 58 cursos avaliados com conceito 5, 36 eram vinculados a Universidades.
Os cursos vinculados a esse tipo de Organização Acadêmica tiveram o conceito 3 como conceito modal, atribuído a 206 cursos ( 41,7 \% ).
Os demais cursos vinculados a Universidades receberam os conceitos 1 ( 17 cursos, 3,4 \% ), 2 ( 80 cursos, 16,2 \% ), 4 ( 147 cursos, 29,8 \% ) e conceito 5 ( 36 cursos, 7,3 \%, como já mencionado ).
Ficaram Sem Conceito ( SC ), oito cursos ( 1,6 \% ).
Entre os cursos em Centros Universitários, o conceito modal foi também 3, atribuído a 63 cursos ( 44,7 \% ).
Dois cursos ( 1,4 \% ) vinculados a esse tipo de Organização Acadêmica ficaram Sem Conceito ( SC ), e os outros cursos receberam os conceitos 2 ( 33 cursos, 23,4 \% ), 4 ( 36 cursos, 25,5 \% ), 5 ( sete cursos, 5,0 \% ), não sendo atribuído conceito 1 a nenhum curso de Centros Universitários.
Dos 570 cursos mantidos por Faculdades, 244 ( 42,8 \% ) receberam conceito 3, o conceito modal.
O segundo conceito mais frequentemente atribuído aos cursos oferecidos por essa categoria de IES foi o conceito 2, recebido por 200 cursos ( 35,1 \% ).
Na sequência, 85 cursos obtiveram conceito 4 ( 14,9 \% ) e 14 cursos ( 2,5 \% ), conceito 5.
Um curso ( 0,2 \% ) ficou Sem Conceito ( SC ) e 26 cursos ( 4,6 \% ) receberam conceito 1.
Dos sete cursos mantidos por CEFET/IFET, seis ( 85,8 \% ) receberam conceitos 3 e 4 ( três em cada ), os conceitos modais.
O outro curso avaliado oferecido por essa categoria de IES recebeu o conceito 5, ( 14,3 \% ).
Nenhum curso ficou Sem Conceito ( SC ) ou recebeu os conceitos 1 e 2.
Considerando -se separadamente as regiões brasileiras, verifica -se que, na região Norte, as Universidades participaram com 45 dos 88 cursos participantes ( 51,1 \% ).
O conceito modal foi o conceito 2, recebido por 18 cursos ( 40,0 \% ).
Os demais cursos receberam os conceitos 1 ( sete cursos, 15,6 \% ), 3 ( 14 cursos, 31,1 \% ) e 4 ( seis cursos, 13,3 \% ).
Nenhum curso ficou Sem Conceito ( SC ) ou recebeu conceito 5 nesta combinação de categoria e região.
Os Centros Universitários da região Norte foram representados por cinco cursos ( 5,7 \% ), aos quais foram atribuídos os conceitos 2 e 3 ( recebido por dois cursos, 40,0 \%, cada ).
Nenhum curso ficou Sem Conceito ( SC ) ou recebeu os conceitos 1 e 4.
Um curso recebeu conceito 5 ( 20,0 \% ) nesta combinação de categoria e região.
As Faculdades participaram com 37 cursos na região Norte ( 42,0 \% ) que receberam os conceitos 2 ( 21 cursos, 56,8 \%, o conceito modal ), 3 ( 12 cursos, 32,4 \% ) e 1 ( quatro cursos, 10,8 \% ).
Nenhum curso ficou Sem Conceito ( SC ) ou recebeu os conceitos 4 e 5.
Os CEFET/IFET participaram com um curso ( 1,1 \% ), que recebeu conceito 4.
Na região Nordeste, as Universidades participaram com 114 dos 228 cursos da Área de Pedagogia oferecidos na região ( 50,0 \% ), dos quais dois ( 1,8 \% ) ficaram Sem Conceito ( SC ).
Nota -se que tais cursos apresentaram uma distribuição com moda no conceito 3, atribuído a 56 cursos ( 49,1 \% ).
Obtiveram os conceitos 1, 2, 4 e 5, respectivamente, sete, 23, 21 e cinco cursos, correspondendo a 6,1 \%, 20,2 \%, 18,4 \% e 4,4 \%, respectivamente, dos cursos dessa categoria na região.
Os Centros Universitários contaram com 12 cursos participante na região Nordeste ( 5,3 \% ).
Esses obtiveram conceitos 2, o modal com cinco cursos ( 41,7 \% ), 3 ( quatro cursos, 33,3 \% ), 4 ( um curso, 8,3 \% ) e 5 ( dois cursos, 16,7 \% ).
Nenhum curso oferecido por Centros Universitários no Nordeste ficou Sem Conceito ( SC ) ou recebeu conceito 1.
As Faculdades foram representadas por 102 cursos na região Nordeste ( 44,7 \% ).
Nota -se que tais cursos apresentaram distribuição modal no conceito 2 com 49 cursos ( 48,0 \% ).
O segundo conceito mais atribuído foi o conceito 3, com 33 cursos ( 32,4 \% ).
Os conceitos 2 e 4 foram atribuídos igualmente a 10 cursos ( 9,8 \% ) cada.
Na região Nordeste, nenhum curso oferecido por Faculdades recebeu conceito 5, tampouco ficou Sem Conceito ( SC ).
Os CEFET/IFET não apresentaram cursos nessa região.
Na região Sudeste, as Universidades concentraram 180 dos 556 cursos de Pedagogia da região ( 32,4 \%, o menor percentual dessa categoria entre as regiões ).
Entre os cursos em Universidades, nessa região, o conceito modal foi 4, atribuído a 68 cursos ( 37,8 \% ).
Os demais cursos receberam os conceitos 1 ( um curso, 0,6 \% ), 2 ( 19 cursos, 10,6 \% ), 3 ( 65 cursos, 36,1 \% ) e 5 ( 24 cursos, 13,3 \% ).
Três cursos ( 1,7 \% ) ficaram Sem Conceito ( SC ).
Os Centros Universitários participaram com 83 cursos na região Sudeste ( 14,9 \%, o maior percentual dessa categoria entre as regiões ), dos quais 39 ( 47,0 \% ) obtiveram o conceito modal 3.
O conceito 2 foi atribuído a 17 cursos ( 20,5 \% ), o conceito 4, a 23 cursos ( 27,7 \% ), o conceito 5, a três cursos ( 3,6 \% ) e nenhum curso recebeu conceito 1.
Um curso ( 1,2 \% ) ficou Sem Conceito ( SC ) nessa combinação de Organização Acadêmica e região.
Já as Faculdades foram representadas por 293 cursos na região Sudeste ( 52,7 \%, o maior O Gráfico 5.3 apresenta a distribuição do Conceito Enade dos cursos de Licenciatura segundo a Organização Acadêmica da IES.
Os cursos em CEFET/IFET ( linha roxa ) percentual dessa categoria entre as regiões ), dos quais um curso ( 0,3 \% ) ficou Sem Conceito ( SC ) e quatro receberam conceito 1 ( 1,4 \% ).
Aos demais cursos foram atribuídos os conceitos 2 ( 93 cursos, 31,7 \% ), 3 ( 140 cursos, 47,8 \%, o conceito modal ), 4 ( 44 cursos, 15,0 \% ) e 5 ( 11 cursos, 3,8 \% ).
Os CEFET/IFET não apresentaram cursos nessa região.
Dos 211 cursos da região Sul, 103 eram de Universidades ( 48,8 \% ), para os quais o conceito modal foi 3, atribuído a 45 cursos ( 43,7 \% ).
Nesse tipo de organização, três cursos ( 2,9 \% ) ficaram Sem Conceito ( SC ) e um curso ( 1,0 \% ) recebeu conceito 1.
Os demais cursos receberam os conceitos 2 ( seis cursos, 5,8 \% ), 4 ( 41 cursos, 39,8 \% ) e 5 ( sete cursos, 6,8 \% ).
Dos 25 cursos participantes de Centros Universitários da região Sul ( 11,8 \% ), nenhum ficou Sem Conceito ( SC ) ou recebeu conceito 1.
O conceito modal, 3, foi atribuído a 12 cursos ( 48,0 \% ).
Os conceitos restantes, 2, 4 e 5, foram atribuídos a quatro, oito e um cursos, respectivamente, correspondendo a 16,0 \%, 32,0 \% e 4,0 \% dos cursos dessa categoria na região.
Quanto aos 79 cursos vinculados a Faculdades na região Sul ( 37,4 \% ), 42 receberam o conceito modal 3 ( 53,2 \% ).
O conceito 2 foi atribuído a 13 cursos ( 16,5 \% ), e o conceito 4, a 21 cursos ( 26,6 \% ).
Nenhum curso ficou Sem Conceito ( SC ) ou recebeu conceito 1.
O conceito 5 foi atribuído a três cursos ( 3,8 \% ).
Nesta região foram avaliados quatro cursos em CEFET/IFET, dois cursos receberam o conceito 3 e dois, o conceito 4.
Na região Centro-Oeste, 52 dos 129 cursos eram de Universidades ( 40,3 \% ).
Nesse tipo de organização, nenhum curso ficou Sem Conceito ( SC ).
O conceito 1 foi atribuído a um curso ( 1,9 \% ), o conceito 2, a 14 cursos ( 26,9 \% ) e o conceito 3, conceito modal, a 26 cursos ( 50,0 \% ).
Além disso, 11 cursos ( 21,2 \% ) alcançaram o conceito 4 e nenhum curso recebeu o conceito máximo, 5.
Os Centros Universitários da região Centro-Oeste contaram com 16 cursos ( 12,4 \% ), que obtiveram os conceitos 3 ( seis cursos, 37,5 \%, o conceito modal ) e 4 ( quatro cursos, 25,0 \% ).
Dos 16 cursos oferecidos em Centros Universitários na região Centro-Oeste, cinco receberam o conceito 2 ( 31,3 \% ).
Nenhum curso recebeu os conceitos 1 ou 5 e apenas um curso ( 6,3 \% ) ficou Sem Conceito ( SC ).
Nesta região foram dois os cursos avaliados em CEFET/IFET, que receberam os conceitos 3 e 5.
Nas Faculdades, que representaram 45,7 \% dos cursos participantes na região Centro-Oeste, o conceito modal foi 2 atribuído a 24 cursos ( 40,7 \% ).
Os conceitos 1, 3 e 4 foram atribuídos, respectivamente, a oito, 17 e dez cursos, correspondendo a 13,6 \%, 28,8 \% e 16,9 \% dos cursos dessa categoria na região.
Nenhum curso ficou Sem Conceito ( SC ) ou recebeu conceito 5.
O Gráfico 5.3 apresenta a distribuição do Conceito Enade dos cursos de Licenciatura segundo a Organização Acadêmica da IES.
Os cursos em CEFET/IFET ( linha roxa ) apresentam uma poligonal mais à direita do que as demais, denotando uma melhor distribuição de Conceitos Enade.
No outro extremo, temos os cursos oferecidos em Faculdades ( linha vermelha ) com a poligonal mais à esquerda e, consequentemente, uma distribuição de Conceitos Enade com menores valores.
IES de Centros Universitários ( linha verde ) e em Universidades ( linha azul ) apresentam uma situação intermediária.

% matched lemmas: distribuir, elevado, enfatizar
\end{textsample}
